%% kapa: exam=1 round=36 year=2025 period=1 form=A subject=real_estate subject_ko=부동산학원론 seq=15 number=95
%% source: Q-Net 감정평가사 기출문제 (cst00309) artlSeq=5231525 file=raw/1008/5231525/2025년 제36회 감정평가사 1차시험 1교시 문제지 원본.pdf
%% conversion: kordoc 4.24.51 -> tools/gen1.py (auto); review=proofread-sonnet; flags=none
\begin{question}{95}
토지공개념에 관한 설명으로 옳지 않은 것은?
\begin{choices}
  \item 토지의 공익성을 강조하는 개념으로 정부가 공공의 이익을 위하여 토지의 소유권을 제한할 수 있다는 인식을 반영하고 있다.
  \item 국가는 국민 모두의 생산 및 생활의 기반이 되는 국토의 효율적이고 균형있는 이용ㆍ개발과 보전을 위하여 법률이 정하는 바에 따라 제한과 의무를 과(課)할 수 있다는 근거를 반영하고 있다.
  \item 토지의 사유재산권을 부정하고 그 보유ㆍ이용ㆍ개발이 공공복리에 적합해야 한다고 보았다.
  \item 지가의 폭등과 개발이익환수의 미비로 인해 만연한 토지투기를 근절하고자 토지공개념을 도입하였다.
  \item 개발부담금제도는 토지로부터 발생되는 개발이익을 환수하여 이를 적정하게 배분하여서 토지에 대한 투기를 방지하고 토지의 효율적인 이용을 촉진하기 위한 제도이다.
\end{choices}
\end{question}
\answer{3}
